\documentclass[12pt,a4paper]{article}

% ─── Packages ────────────────────────────────────────────────────────────────
\usepackage[utf8]{inputenc}
\usepackage[T1]{fontenc}
\usepackage{geometry}
\geometry{margin=2.5cm}

\usepackage{amsmath,amssymb}
\usepackage{booktabs}
\usepackage{array}
\usepackage{xcolor}
\usepackage{enumitem}
\usepackage{tcolorbox}
\tcbuselibrary{skins, breakable}
\usepackage{hyperref}

% ─── Colors ──────────────────────────────────────────────────────────────────
\definecolor{urgence}{RGB}{180,30,30}
\definecolor{theorie}{RGB}{30,80,160}
\definecolor{pratique}{RGB}{20,130,80}
\definecolor{attention}{RGB}{200,120,0}
\definecolor{bleuléger}{RGB}{230,240,255}
\definecolor{vertléger}{RGB}{230,248,238}
\definecolor{rougeclair}{RGB}{255,235,235}
\definecolor{orangeclair}{RGB}{255,245,225}
\definecolor{bleupolicier}{RGB}{225,235,255}

% ─── Custom boxes ─────────────────────────────────────────────────────────────
\newtcolorbox{vehiclebox}[2]{
  colback=#1, colframe=#2, fonttitle=\bfseries\large,
  title={##1}, breakable, sharp corners=south,
  left=8pt, right=8pt, top=6pt, bottom=6pt
}

\newtcolorbox{synthbox}{
  colback=gray!6, colframe=gray!40, fonttitle=\bfseries,
  title={Synthèse : influence sur le choix d'itinéraire},
  breakable, sharp corners=south
}

% ─── Hyperref ─────────────────────────────────────────────────────────────────
\hypersetup{
  colorlinks=true,
  linkcolor=theorie,
  urlcolor=urgence,
  pdftitle={Priorités et contraintes des véhicules d'urgence},
  pdfauthor={Songo — Simulation de routage},
}

% ─────────────────────────────────────────────────────────────────────────────
\begin{document}

\begin{center}
  {\LARGE\bfseries Priorités et contraintes des véhicules d'urgence}\\[6pt]
  {\large Simulation de routage Songo — Yaoundé}\\[3pt]
  {\small\color{gray} MVP · Routage dynamique A*}
\end{center}

\vspace{0.5cm}

\noindent
Dans le MVP de simulation, trois types de véhicules d'urgence sont disponibles :
l'ambulance, le camion de pompiers et le véhicule de police.
Bien qu'ils partagent le même moteur de routage (algorithme A* avec poids dynamiques),
\textbf{leurs contraintes et leurs priorités diffèrent}, ce qui conduit le système
à calculer des itinéraires différents à départ et destination identiques.

\bigskip

% ── Ambulance ────────────────────────────────────────────────────────────────
\begin{tcolorbox}[colback=rougeclair, colframe=urgence, fonttitle=\bfseries\large,
  title={[Ambulance] Ambulance — Priorité absolue : rapidité},
  breakable, sharp corners=south]

L'ambulance opère sous la contrainte la plus forte : \textbf{chaque seconde compte}.
Le moteur lui attribue un coefficient de priorité maximal sur le temps de trajet,
et minimise avant tout la durée, quitte à accepter des routes plus difficiles.

\begin{itemize}[leftmargin=1.5em, itemsep=3pt]
  \item \textbf{Axes étroits autorisés} : les rues secondaires ou les ruelles sont
        empruntables si elles raccourcissent le trajet ; la largeur de la voie
        ne constitue pas un filtre bloquant.
  \item \textbf{Routes dégradées tolérées} : une piste en terre sera retenue si
        elle évite un embouteillage sur un axe principal.
  \item \textbf{Pénalité pluie modérée} : même sous la pluie, le moteur maintient
        une partie des itinéraires rapides, quitte à augmenter légèrement le risque.
  \item \textbf{Nuit} : aucune restriction supplémentaire ; la priorité temporelle
        prime sur toute contrainte de circulation nocturne.
\end{itemize}

\smallskip
\textit{Résultat dans la simulation :} l'ambulance obtient systématiquement
le chemin le plus court en temps, même s'il est sinueux ou peu praticable
pour un véhicule plus lourd.

\end{tcolorbox}

\bigskip

% ── Pompiers ─────────────────────────────────────────────────────────────────
\begin{tcolorbox}[colback=orangeclair, colframe=attention, fonttitle=\bfseries\large,
  title={[Pompiers] Camion de pompiers — Contraintes physiques fortes},
  breakable, sharp corners=south]

Le camion de pompiers est un \textbf{véhicule lourd} : son gabarit et son poids
imposent des contraintes physiques que le moteur doit respecter en priorité,
avant même d'optimiser la durée.

\begin{itemize}[leftmargin=1.5em, itemsep=3pt]
  \item \textbf{Routes non goudronnées exclues} : les pistes en terre ou les voies
        non stabilisées reçoivent un poids prohibitif (ou sont bloquées) ;
        le moteur les contourne systématiquement.
  \item \textbf{Voies trop étroites exclues} : les ruelles inférieures à un gabarit
        minimal ne sont pas empruntables, même si elles sont plus rapides.
  \item \textbf{Zones civiles interdites accessibles} : en contrepartie, le camion
        peut traverser certains axes normalement fermés aux civils
        (zones réservées, cours intérieures de bâtiments publics, etc.).
  \item \textbf{Pénalité pluie élevée} : sur sol détrempé, les routes non goudronnées
        sont totalement éliminées, resserrant encore l'espace des routes valides.
\end{itemize}

\smallskip
\textit{Résultat dans la simulation :} le camion emprunte presque toujours les
grands axes goudronnés, même si cela allonge le trajet ; il peut toutefois
se permettre certains raccourcis interdits aux véhicules civils.

\end{tcolorbox}

\bigskip

% ── Police ───────────────────────────────────────────────────────────────────
\begin{tcolorbox}[colback=bleupolicier, colframe=theorie, fonttitle=\bfseries\large,
  title={[Police] Police — Équilibre rapidité / accessibilité},
  breakable, sharp corners=south]

Le véhicule de police adopte un \textbf{profil intermédiaire} : il cherche un
compromis entre la rapidité et la faisabilité du trajet, sans la contrainte
physique du poids ni l'urgence vitale de l'ambulance.

\begin{itemize}[leftmargin=1.5em, itemsep=3pt]
  \item \textbf{Axes étroits partiellement autorisés} : les voies secondaires sont
        acceptées si leur état est correct, mais le moteur leur applique un léger
        malus par rapport aux grands axes.
  \item \textbf{Routes dégradées avec malus} : une piste en terre n'est pas exclue
        mais reçoit un coefficient défavorable ; elle ne sera choisie qu'en
        dernier recours.
  \item \textbf{Priorité tempérée} : contrairement à l'ambulance, la police ne
        prend pas de risques excessifs pour gagner quelques secondes.
  \item \textbf{Nuit} : certains axes nocturnes peuvent être évités pour des
        raisons de sécurité opérationnelle, sans que cela bloque complètement
        l'itinéraire.
\end{itemize}

\smallskip
\textit{Résultat dans la simulation :} la police obtient un trajet souvent proche
de celui de l'ambulance, mais légèrement plus long — le moteur préfère un chemin
sûr et accessible à un raccourci risqué.

\end{tcolorbox}

\bigskip

% ── Synthèse ─────────────────────────────────────────────────────────────────
\begin{synthbox}
\noindent
Le tableau suivant résume les différences de traitement dans le moteur A* :

\medskip
\begin{tabular}{lccc}
  \toprule
  \textbf{Critère} & \textbf{Ambulance} & \textbf{Pompiers} & \textbf{Police} \\
  \midrule
  Objectif principal     & Temps minimal    & Gabarit / sol    & Équilibre \\
  Axes étroits           & \checkmark autorisés      & $\times$ bloqués        & $\sim$ malus   \\
  Routes non goudronnées & $\sim$ tolérées       & $\times$ bloquées       & $\sim$ malus   \\
  Zones civiles fermées  & $\sim$ selon contexte & \checkmark autorisées     & $\times$ non     \\
  Pénalité pluie         & Faible           & Très élevée      & Modérée   \\
  Contrainte nocturne    & Aucune           & Aucune           & Légère    \\
  \bottomrule
\end{tabular}

\medskip
\noindent
Ces paramètres sont encodés comme \textbf{poids dynamiques} dans le graphe routier.
Changer de type de véhicule avant de lancer la simulation suffit à modifier
l'ensemble des coefficients, et donc l'itinéraire proposé — même pour un trajet
départ–destination identique.
\end{synthbox}

\end{document}